% TOP_PROBLEM: {"id":30007631,"problem_number":"LOCAL-30007631","title":"Contracting frictions in global sourcing","category_id":21,"category":{"id":21,"name":"economics","display_name":"Economics","description":"Open economic theory statements, published research agendas, and proposed scoped specifications; question provenance and historical priority are qualified per record.","slug":"economics","order_index":21,"created_at":"2026-10-08T00:00:00Z"},"status":"provenance_pending","external_url":null,"legacy_ids":[],"published":false,"statement":"In the explicitly specified setting: How much do contract enforcement and organizational frictions limit the gains from global sourcing and multinational production? The mathematical statement above fixes the target, assumptions and scope of this catalogue formulation.","economics":{"original_id":120,"field":"Trade, globalization and geoeconomics","kind":"identification","formalization_basis":"proposed_specification","source_status":"not_verified","priority_status":"not_established","priority_note":"No primary passage stating this exact catalogue question as open was verified; no original-priority claim is made.","earliest_verified_reference":null,"current_reference":{"authors":["Davin Chor","Lin Ma"],"title":"The Aggregate Welfare and Trade Implications of Contracting Frictions in Global Sourcing","year":2025,"url":"https://lin-ma.com/uploads/3/6/0/7/36070314/w34044.pdf","doi":"10.3386/w34044","locator":"Abstract and introduction","evidence_summary":"Provides a particular model and estimates, not verified open statement."},"formalization_note":"The cited current paper constructs and estimates a contracting model. Its results solve conditional subcases; the proposed broader identification task is not attributed as an original source statement."},"statusQualification":"Provenance pending: an explicit published statement of this exact open question was not verified. This is a catalogue-authored candidate specification, not a certified open conjecture. The cited current paper constructs and estimates a contracting model. Its results solve conditional subcases; the proposed broader identification task is not attributed as an original source statement.","statusReviewedAt":"2026-10-08"}
% FIRST_POSTED: 2026-10-08
% ENTRY_KIND: statement_only
% !TeX program = lualatex
\documentclass[11pt]{article}
\usepackage{fontspec}
\usepackage[a4paper,margin=25mm]{geometry}
\usepackage{amsmath,amssymb,mathtools}
\usepackage{xurl}
\usepackage[hidelinks]{hyperref}
\newcommand{\sref}[2]{\hyperlink{source:#1:#2}{[S#2]}}
\setlength{\emergencystretch}{3em}
\begin{document}
% problemId:30007631
\section*{Contracting frictions in global sourcing}\label{30007631}
\subsection{Definitions and mathematical statement}
Fix heterogeneous firms choosing domestic or foreign suppliers and whether to own production stages. Contracts are incomplete in a specified sense: a declared subset of investment or quality actions cannot be enforced by courts. Let \(\kappa_{ab}\) measure enforcement costs or noncontractible exposure for transactions between jurisdictions \(a,b\). Firms choose organization, sourcing and relationship-specific investment optimally, with wages and prices clearing. Let \(W(\kappa)\) be discounted household welfare and \(Q(\kappa)\) trade volumes. The task is to identify losses from these frictions relative to a defined counterfactual \(\kappa=0\):
\[\Delta W=W(0)-W(\kappa),\qquad\Delta Q=Q(0)-Q(\kappa).\]
Specify which institutional change implements friction removal and which technologies, preferences and trade policies remain fixed. Separate enforcement frictions from search, transport and financing costs using a stated exclusion or structural restriction. Establish whether organizational choices and sourcing patterns uniquely identify \(\kappa\) or only composite wedges. Include endogenous multinational ownership, entry and supplier investment in welfare comparisons. Report sensitivity to bargaining rules and cross-country institutional measurement. A model matching observed sourcing shares can estimate a conditional counterfactual without proving that its enforcement mechanism is uniquely identified or that the broader research question was first stated in that paper.

\par\textbf{Formulation and scope.} The cited current paper constructs and estimates a contracting model. Its results solve conditional subcases; the proposed broader identification task is not attributed as an original source statement.

\par\textbf{Open-question provenance.} An explicit published open-question statement was not verified. Sources below are supporting literature and must not be treated as a first listing.

\par\textbf{Historical priority.} No primary passage stating this exact catalogue question as open was verified; no original-priority claim is made.

\subsection{Short English statement}
In the explicitly specified setting: How much do contract enforcement and organizational frictions limit the gains from global sourcing and multinational production? The mathematical statement above fixes the target, assumptions and scope of this catalogue formulation.

\subsection{Sources}
\begin{itemize}\raggedright
\item[\textnormal{[S1]}]\hypertarget{source:30007631:1}{} Davin Chor, Lin Ma. 2025. The Aggregate Welfare and Trade Implications of Contracting Frictions in Global Sourcing. Abstract and introduction.
\url{https://lin-ma.com/uploads/3/6/0/7/36070314/w34044.pdf}.
DOI: 10.3386/w34044.
{\small\textit{Supporting or current literature; see evidence qualification. Provides a particular model and estimates, not verified open statement.}}
\end{itemize}
\end{document}
